% This must be in the first 5 lines to tell arXiv to use pdfLaTeX, which is strongly recommended.
\pdfoutput=1
% In particular, the hyperref package requires pdfLaTeX in order to break URLs across lines.

\documentclass[11pt]{article}

% Remove the "review" option to generate the final version.
\usepackage{acl}

% Standard package includes
\usepackage{times}
\usepackage{latexsym}

% For proper rendering and hyphenation of words containing Latin characters (including in bib files)
\usepackage[T1]{fontenc}
% For Vietnamese characters
% \usepackage[T5]{fontenc}
% See https://www.latex-project.org/help/documentation/encguide.pdf for other character sets

% This assumes your files are encoded as UTF8
\usepackage[utf8]{inputenc}

% This is not strictly necessary, and may be commented out.
% However, it will improve the layout of the manuscript,
% and will typically save some space.
\usepackage{microtype}

% This is also not strictly necessary, and may be commented out.
% However, it will improve the aesthetics of text in
% the typewriter font.
\usepackage{inconsolata}


% If the title and author information does not fit in the area allocated, uncomment the following
%
%\setlength\titlebox{<dim>}
%
% and set <dim> to something 5cm or larger.

\title{Character Classification in Game of Thrones Using Language Models Trained on Screenplay Dialogue}

% Author information can be set in various styles:
% For several authors from the same institution:
% \author{Author 1 \and ... \and Author n \\
%         Address line \\ ... \\ Address line}
% if the names do not fit well on one line use
%         Author 1 \\ {\bf Author 2} \\ ... \\ {\bf Author n} \\
% For authors from different institutions:
% \author{Author 1 \\ Address line \\  ... \\ Address line
%         \And  ... \And
%         Author n \\ Address line \\ ... \\ Address line}
% To start a seperate ``row'' of authors use \AND, as in
% \author{Author 1 \\ Address line \\  ... \\ Address line
%         \AND
%         Author 2 \\ Address line \\ ... \\ Address line \And
%         Author 3 \\ Address line \\ ... \\ Address line}

\author{John Brassey \\
  \texttt{jabrassey@gmail.com}
        }

\begin{document}
\maketitle
\begin{abstract}
This project sought to answer the question of whether a relatively small language model could be fine tuned to classify, by character name, lines of dialogue from the television series Game of Thrones. The larger aim of the project was to contribute in a small way to the question of whether language models can generate innovative and creative fiction content. The author hypothesized that if a language model could be trained to identify characters by their defining trait (what they say), it could eventually be trained to generate unique "voices" for characters in a story. Based on the results of this project, small language models may struggle to classify lines of dialogue by character name in a series with large numbers of main characters.
\end{abstract}

\section{Introduction}

Although large language models (LLMs) were first "transformed" by the publication of the now-famous "Attention is All You Need" paper \cite{vaswani2017attention} (and its corresponding transformer/attention architecture) in 2017, it wasn't until the introduction of ChatGPT to the public in November of 2022 that LLMs, and their capabilities,  became the subject of popular conversation. Suddenly, everyone from high school teenagers, to politicians, to screenwriters was talking about the potential, and potential threat, of language models like ChatGPT. Would they make software developers obsolete? Would creative writers be out of a job? Questions that had only ever shown up in science fiction novels now took on global importance.

Having received my first degree in English Literature, I was curious about the quality of creative fiction that an LLM like ChatGPT could produce. Writing great fiction is hard, and creating unique, believable, sympathetic characters is even harder. Given that some studies have already suggested that LLMs score worse on metrics of creativity than humans \cite{chakrabarty2023art}, I wanted to start with a simple research question: can an LLM differentiate between different characters in a television show based on their lines of dialogue? Before a person or a language model can be expected to create a set of unique characters for a story (and give them each their own individual way of speaking), it seems reasonable that they (or it) should be able to identify characters based solely on the way that they speak. That is to say, if the screenwriters have done their job, each character in a show should be identifiable using only their on-screen dialogue.

Game of Thrones is a show renowned for its diverse cast of interesting characters, each with their own backgrounds, desires, and machinations, and so it seemed like the perfect show on which to test a language model's ability to classify fictional characters by the way they speak.

\section{Background}

Several studies have already attempted to address the question of whether or not language models are capable of genuine creativity. In their 2024 paper titled "\textit{Art of Artifice}? Large Language Models and the False Promise of Creativity," Tuhin Chakrabarty et al. \cite{chakrabarty2023art} argue that language models perform worse than humans on tasks involving creativity. Building on the widely accepted Torrance Test of Creative Thinking (TCTT), the authors, in collaboration with creative writing experts, developed the Torrance Test of Creative Writing (TTCW) to objectively measure what they termed "creativity as a product." In their paper, they showed that, although certain LLMs can perform moderately well on certain creative tasks, assessed holistically, humans significantly outperform language models at creative writing tasks.

Other studies have come to similar conclusions. Giorgio Franceschelli and Mirco Musolesi, in their paper "On the Creativity of Large Language Models" \cite{franceschelli2023creativity} also talk about the concept of a creativity "product," using Margaret Boden's definition of creativity as "the ability to come up with ideas or artifacts that are \textit{new}, \textit{surprising}, and \textit{valuable}" \cite{boden2003autodidacts}. Ultimately, Franceschelli and Musolesi conclude that while language models can at times exhibit some novelty, they are not capable of "transformational creativity."

With several studies having already concluded that LLM's are not currently capable of "transformational creativity" in the realm of creative writing, the goal of this project was to see if LLMs could at least take the first step. A person generally learns to read before they learn to write, and if an LLM is ever going to be capable of generating truly unique and creative works of fiction, it must first be able to identify characters based on their most defining quality: their dialogue.

\section{Data}

The data for this project came from a csv file containing all spoken dialogue in the popular Game of Thrones series, which aired on HBO from 2011 - 2019. The file contained dialogue data by episode, including the original episode air date, the season number, episode number, the character who spoke each line of dialogue, as well as the dialogue itself. In order to test whether an LLM can identify a character solely by their spoken dialogue, this project made use of only the character name and dialogue features. It would have been possible to use other features from this dataset in order to help the model better classify characters. For example, episode number and season number could have been valuable, as certain characters appeared more frequently in certain seasons and within specific episodes within a season, and some characters died throughout the course of the series, meaning that the model, if trained on such data, would not have had to consider classifying certain characters beyond a specific episode in the series. However, in the interests of focusing on a narrowly defined task, only the two columns stated above were used.

Additionally, I derived two features from those original two columns: sentence length (where "sentence" refers to each line of dialogue) and the data type of the "sentence." This last feature was necessary in order to deal with an error where a sentence, classified as a float, was causing errors in the code. With the data prepared and the features chosen, model preparation and training began.

\section{Model and Approach}

As stated in the introduction, the research question of this project was whether a language model, fine tuned  on a small language model using the Game of Thrones dataset described in the Data section of this paper, could classify characters based solely on their lines of dialogue. The high-level approach I took was to conduct a series of experiments to determine which of a plethora of different models performed best, and then to optimize the final model as much as possible.

Given a show with as many characters as Game of Thrones, particularly a show where characters die frequently and often speak only a few lines of dialogue, the first challenge was to determine how much data to train the model on. A brief analysis of the dataset showed that the range for the number of lines of dialogue spoken by a character went from a maximum of 1760 lines of dialogue for one character, to one line of dialogue for 141 characters. This spread necessitated a reduction in the number of characters, as no model could be expected to correctly identify a character who only spoke one line of dialogue throughout the entire series, if for no other reason than that it would be impossible for any characters who only spoke one line of dialogue to show up in both the training and test sets. Thus, the first set of experiments involved reducing the number of characters in the data set.

I created several experiments, each of which further reduced the total number of characters in the dataset. Experiment 1 only included characters who spoke at least one hundred lines of dialogue throughout the series, Experiment 2 only included characters who spoke at least 500 lines of dialogue, and Experiment 3 only included characters who spoke at least 700 lines of dialogue.

Each experiment also contained several sub experiments, which filtered the data based on the number of words in a sentence. Given that the whole point of the project was to build a language model to identify characters based on their unique ways of speaking, asking a model to identify the character associated with a line of dialogue containing only a few words seemed unreasonable. Thus, each of the three original experiments contained three sub experiments, a, b, and c, which filtered the dataset to only include rows where the line of dialogue was longer than 10 "words," 20 "words," or 40 "words." Words is put in quotations here to indicate that sentences were split on white space and thus could include tokens such as punctuation marks.

The sub experiments were altered starting with Experiment 2 to include sentences with lengths greater than or equal to 10, 15, and 20 words, respectively, after it was found during Experiment 1 that only including sentences with lengths greater than 40 words not only shrunk the dataset too significantly, but also resulted in poor model performance.

The next round of experiments, after identifying the best combination of total characters and sentence lengths, consisted of testing different models on which to fine tune the data set. Experiment 4 fine tuned using the RoBERTa model \cite{liu2019roberta},  and Experiment 5 fine tuned using the XLNet model \cite{yang2019xlnet}.

Finally, after the best model (BERT) \cite{devlin2018bert} was identified for fine tuning, the hyper parameters (batch size, max length, number of epochs, and learning rate) were optimized to further increase the performance of the final model.

The idea behind reducing the number of characters was that a model would be much more likely to be able to differentiate between a smaller number of characters, who each had many lines of dialogue and thus were given many opportunities to establish their own unique way of speaking, than it would with a correspondingly larger number of characters, many of whom only spoke a few lines of dialogue. Additionally, as already stated above, the dataset was reduced based on sentence length in order to give the model long enough sentences to be able to realistically identify who spoke them. These initial hypotheses were mostly born out during the evaluation portion of the project.

\section{Evaluation}

\subsection{Task and Procedure}

Evaluation took place throughout the project, and each individual model was trained, tested, and evaluated as it was created.

At the start of each experiment, the dataset was reduced by filtering the characters in the dataset based on how many lines of dialogue they spoke throughout the series. Then, the dataset was further refined by only including samples whose sentences met or exceeded the specified length for the experiment. Once the reduced dataset was obtained for each experiment, the baselines were established, the labels (character names) were encoded, and the dataset was split into training and testing sets using a random seed for reproducibility. The tokenizer was then instantiated using the first sample in the training dataset and the chosen model for the experiment.

Next, the conditions for training and testing were set up. That is to say, the hyper parameters were set for both training and testing, the samples were tokenized, and the training and testing loaders were created. Then for each model, a loop consisting of the specified number of epochs was run, during which the training loss and test accuracy were calculated.

\subsection{Metrics}

The metric used for evaluation of each model in this project was a simple accuracy score. Since the chosen task was a classification task, the easiest and most effective way to evaluate each model's performance was simply to divide the number of samples it correctly predicted by the size of the given dataset. Additionally, two baselines were calculated for each experiment: a random baseline, and a most common baseline. In all cases, the most common baseline resulted in a larger value, so it was always the baseline against which each model's performance was compared.

\subsection{Results}

I initially trained a model on all the data, but the accuracy was so poor (below 6\%) that I immediately designed different experiments to test the efficacy of reducing the dataset based on number of lines of dialogue and sentence length.

Experiment 1 consisted of only characters who spoke at least 100 lines of dialogue throughout the show. Sub experiments a, b, and c further reduced the dataset by only choosing samples which had sentence lengths greater than 10, 20, and 40, respectively. Sub experiment a, with a most common baseline of 12.1\%, achieved a test accuracy of 12.35\%. Sub experiment b, with a most common baseline of 13.4\%, achieved a test accuracy of 16.52\%. Sub experiment c, with a most common baseline of 14.8\%, achieved a test accuracy of 15.23\%.

Experiment 2 consisted of only characters who spoke at least 500 lines of dialogue throughout the series. Experiment 2 was only tested with sentence lengths of 20 "words" or more. With a most common baseline of 28.9\%, Experiment 2 achieved a test accuracy of 32.97\%.

Experiment 3, which only included characters who spoke at least 700 lines of dialogue throughout the series, contained two sub experiments, a and b, which used samples that contained sentence lengths of at least 15 and 20 "words,", respectively. Sub experiment a, with a most common baseline of 30.8\%, achieved a test accuracy of 42.88\%. Sub experiment b, with a most common baseline of 31.12\%, achieved a test accuracy of 38.58\%.

Since Sub experiment 3a performed best during the initial series of tests, the experiments using different models for fine tuning were only tested using samples whose characters spoke at least 700 lines of dialogue and whose sentences were 20 "words" or longer.

Experiment 4 used the RoBERTa model for fine tuning. with a most common baseline of 30.8\%, Experiment 4 achieved a test accuracy of 32.15\%. Experiment 5, which used the XLNet model for fine tuning, achieved a test accuracy of 32.15\% while maintaining the same most common baseline as Experiment 4.

Experiment 3a performed the best of all models tested up to that point, so hyper parameter optimization took place on that model.

Experiment 6 ran a loop to find the best hyper parameters (max length, batch size, number of epochs, and learning rate), but none were found that improved upon the initial parameters used in Experiment 3a.

Finally, a final, optimized model was run for 30 epochs. Against a most common baseline of 30.8\%, the model achieved an accuracy on the test set of 43.29\%

\section{Implications}

The implications of this project, are, I believe, relatively limited in scope. Given the black box nature of many language models, it is difficult to tell why the accuracy of the model was so low. Although it outperformed both baselines, I would not suggest using it as a reliable classifier for this dataset. However, that fact alone may have its own implications. The fact that the model had difficulty classifying characters by their spoken lines of dialogue may give further evidence to the claim that language models cannot generate creative works of fiction, and thus this project may be useful to researches interested in the ability of LLMs to produce novel and creative fiction content. 

Given that this project was simply a classification task for characters on a television show, the author of this paper does not anticipate it having any negative effects on other people. In fact, this project may prove reassuring to those who work in the creative writing industry professionally, as this project gives some evidence that LLMs of the scope of BERT and RoBERTa do a moderate-to-poor job of identifying characters based solely on their spoken dialogue.

\section{Conclusion}

There are two potential conclusions to be drawn from this project. The first is that the writers for Game of Thrones, and, by extension, George R.R. Martin, did not do a good enough job of making characters distinguishable based solely on their dialogue.

The other possible conclusion is that small LLMs such as BERT and RoBERTa have room for improvement when it comes to identifying characters in a story based solely on lines of dialogue. Extrapolating the results of this paper even further, it may also be the case, though this paper is a long way from proving such a claim, that small LLMs have a ways to go before they are capable (if ever) of producing truly original, creative fiction content, since they struggled in this instance to even identify characters based on their defining feature: what they say.

\section{Future Improvements}

Going forward, the author would like to take the results of this project even further. He would eventually like to create a dialogue generating model that, when prompted with the name of a major character from the Game of Thrones television series ("major character" in this instance being a character who speaks more than 700 lines of dialogue over the course of the series), will generate text in the "voice" of that character. The author would also like to create a simple website with a user interface allowing the user to input the character of their choice and to see output of model-generated dialogue in the "voice" of that character. However, given the relatively poor results of the character classification model, the author is now skeptical that such a dialogue generation model would even function properly.

\bibliographystyle{acl_natbib}
\bibliography{custom}

\end{document}
